\documentclass[aspectratio=169,11pt]{beamer}
\usetheme{Madrid}
\usecolortheme{dolphin}
\setbeamertemplate{navigation symbols}{}
\setbeamertemplate{caption}[numbered]

\usepackage[utf8]{inputenc}
\usepackage[T1]{fontenc}
\usepackage[spanish,es-noshorthands]{babel}
\usepackage{amsmath,amssymb,amsthm,mathtools}
\usepackage{tikz}
\usepackage{pgfplots}
\pgfplotsset{compat=1.18}
\usepackage{booktabs}
\usepackage{array}

\title[Prob. y Estadística]{Clase 5: Definición de variable aleatoria y función de distribución}
\subtitle{Unidad 2: Variables aleatorias}
\institute[UNS]{Universidad Nacional del Sur \\ Departamento de Matemática}
\author{Probabilidad y Estadística (5765)}
\date{}

\begin{document}

\begin{frame}
\titlepage
\end{frame}

\begin{frame}{Contenidos}
\tableofcontents
\end{frame}

\section{Motivación}

\begin{frame}{¿Por qué necesitamos variables aleatorias?}
\begin{block}{El problema}
Hasta ahora trabajamos con espacios muestrales $\Omega$ y eventos $A \subseteq \Omega$ arbitrarios (caras, colores, defectuoso/no defectuoso, etc.). Muchas veces nos interesa un \textbf{valor numérico} asociado al resultado del experimento.
\end{block}

\begin{example}
\begin{itemize}
\item Tirar tres monedas: interesa el \emph{número de caras}, no la secuencia exacta.
\item Inspeccionar un lote: interesa la \emph{cantidad de artículos defectuosos}.
\item Medir el tiempo de vida de una lámpara: interesa el \emph{número de horas} hasta que se quema.
\end{itemize}
\end{example}

\begin{alertblock}{Idea clave}
Una \textbf{variable aleatoria} es una función que asigna un número real a cada resultado del experimento aleatorio.
\end{alertblock}
\end{frame}

\section{Definición formal}

\begin{frame}{Variable aleatoria: definición}
\begin{definition}[Variable aleatoria]
Sea $(\Omega, \mathcal{A}, P)$ un espacio de probabilidad. Una \textbf{variable aleatoria} (v.a.) es una función
$$X : \Omega \longrightarrow \mathbb{R}$$
tal que para todo $x \in \mathbb{R}$, el conjunto
$$\{\omega \in \Omega : X(\omega) \le x\}$$
es un evento (pertenece a $\mathcal{A}$), es decir, es \textbf{medible}.
\end{definition}

\begin{alertblock}{Notación}
Escribimos $X$ para la variable aleatoria (la función) y $x$ para un valor particular que puede tomar. El evento $\{\omega : X(\omega) \le x\}$ se abrevia $\{X \le x\}$.
\end{alertblock}
\end{frame}

\begin{frame}{Ejemplo introductorio}
\begin{example}
Se lanzan tres monedas equilibradas. $\Omega = \{CCC, CCX, CXC, XCC, CXX, XCX, XXC, XXX\}$, cada resultado con probabilidad $1/8$.

Definimos $X(\omega) = $ número de caras en $\omega$.

\begin{center}
\begin{tabular}{lccccc}
\toprule
$\omega$ & $XXX$ & $CXX,XCX,XXC$ & $CCX,CXC,XCC$ & $CCC$ \\
\midrule
$X(\omega)$ & 0 & 1 & 2 & 3 \\
\bottomrule
\end{tabular}
\end{center}

Así, $X$ toma valores en $\{0,1,2,3\}$ y por ejemplo
$$P(X=2) = P(\{CCX,CXC,XCC\}) = \frac{3}{8}.$$
\end{example}
\end{frame}

\section{Función de distribución acumulada}

\begin{frame}{Definición de $F(x)$}
\begin{definition}[Función de distribución acumulada]
La \textbf{función de distribución acumulada} (f.d.a.) de una variable aleatoria $X$ es la función $F : \mathbb{R} \to [0,1]$ dada por
$$F(x) = P(X \le x), \qquad x \in \mathbb{R}.$$
\end{definition}

\begin{block}{Interpretación}
$F(x)$ es la probabilidad acumulada de que $X$ tome un valor menor o igual a $x$. Contiene \emph{toda} la información probabilística de $X$.
\end{block}

\begin{example}
Para el ejemplo de las tres monedas:
$$F(1.5) = P(X \le 1.5) = P(X=0) + P(X=1) = \frac{1}{8} + \frac{3}{8} = \frac{4}{8} = \frac{1}{2}.$$
\end{example}
\end{frame}

\begin{frame}{Propiedades de $F(x)$}
\begin{theorem}[Propiedades de la función de distribución]
Toda f.d.a. $F$ satisface:
\begin{enumerate}
\item $F$ es \textbf{monótona no decreciente}: si $x_1 < x_2$ entonces $F(x_1) \le F(x_2)$.
\item $\displaystyle\lim_{x \to -\infty} F(x) = 0$ \quad y \quad $\displaystyle\lim_{x \to +\infty} F(x) = 1$.
\item $F$ es \textbf{continua a derecha}: $\displaystyle\lim_{h \to 0^+} F(x+h) = F(x)$ para todo $x$.
\item $0 \le F(x) \le 1$ para todo $x \in \mathbb{R}$.
\end{enumerate}
\end{theorem}

\begin{alertblock}{Recíproco}
Cualquier función que cumpla (1)--(3) es la f.d.a. de alguna variable aleatoria. Estas propiedades \textbf{caracterizan} a las funciones de distribución.
\end{alertblock}
\end{frame}

\begin{frame}{Demostración (monotonía)}
\begin{proof}[Demostración de la monotonía]
Sean $x_1 < x_2$. Entonces
$$\{X \le x_1\} \subseteq \{X \le x_2\},$$
pues si $X(\omega) \le x_1$ y $x_1 < x_2$, entonces $X(\omega) \le x_2$.

Por la monotonía de la probabilidad (si $A \subseteq B$ entonces $P(A) \le P(B)$):
$$F(x_1) = P(X \le x_1) \le P(X \le x_2) = F(x_2). \qedhere$$
\end{proof}

\begin{block}{Idea de la continuidad a derecha}
Se debe a que trabajamos con $\{X \le x\}$ (desigualdad no estricta). Si tomamos $h_n \downarrow 0$, los eventos $\{X \le x + h_n\}$ decrecen hacia $\{X \le x\}$, y por continuidad de la probabilidad, $F(x+h_n) \to F(x)$.
\end{block}
\end{frame}

\section{Cálculo de probabilidades con \texorpdfstring{$F$}{F}}

\begin{frame}{Probabilidades a partir de $F(x)$}
Usando la f.d.a. podemos calcular la probabilidad de cualquier intervalo:
\begin{block}{Fórmulas básicas}
\begin{align*}
P(X > x) &= 1 - F(x) \\
P(a < X \le b) &= F(b) - F(a), \qquad a<b \\
P(X < x) &= F(x^-) = \lim_{h \to 0^+} F(x-h) \\
P(X = x) &= F(x) - F(x^-) \quad \text{(salto de $F$ en $x$)}
\end{align*}
\end{block}

\begin{alertblock}{Cuidado con los extremos}
$P(a \le X \le b) = F(b) - F(a^-) = F(b) - F(a) + P(X=a)$. Si $X$ es continua, $P(X=a)=0$ y no importa si los intervalos son abiertos o cerrados.
\end{alertblock}
\end{frame}

\begin{frame}{Deducción de $P(a<X\le b) = F(b)-F(a)$}
\begin{proof}
Para $a<b$, escribimos el evento $\{X \le b\}$ como unión disjunta:
$$\{X \le b\} = \{X \le a\} \cup \{a < X \le b\}.$$
Por aditividad de $P$ sobre eventos disjuntos:
$$F(b) = P(X\le b) = P(X \le a) + P(a<X\le b) = F(a) + P(a<X\le b).$$
Despejando: $P(a<X\le b) = F(b) - F(a)$. \qedhere
\end{proof}
\end{frame}

\section{Clasificación de variables aleatorias}

\begin{frame}{Discretas, continuas y mixtas}
\begin{definition}[Variable aleatoria discreta]
$X$ es \textbf{discreta} si toma valores en un conjunto finito o infinito numerable $\{x_1, x_2, \dots\}$. Su f.d.a. $F$ es una función \textbf{escalonada} (constante por tramos, con saltos en los $x_i$).
\end{definition}

\begin{definition}[Variable aleatoria continua]
$X$ es \textbf{continua} si su f.d.a. $F$ es una función \textbf{continua} en todo $\mathbb{R}$ (no tiene saltos), y además existe una función $f \ge 0$ tal que $F(x) = \int_{-\infty}^x f(t)\,dt$.
\end{definition}

\begin{block}{Variables mixtas}
Existen variables \textbf{mixtas}, cuya f.d.a. tiene tramos continuos y también saltos (por ejemplo, tiempos de espera con posibilidad de espera nula). No las estudiaremos en profundidad en este curso.
\end{block}
\end{frame}

\begin{frame}{Gráfico de $F(x)$ para una v.a. discreta}
\begin{example}
Retomando $X=$ número de caras en 3 monedas, con $P(X=0)=1/8$, $P(X=1)=3/8$, $P(X=2)=3/8$, $P(X=3)=1/8$:
\end{example}
\begin{center}
\begin{tikzpicture}[scale=0.85]
\begin{axis}[
  xmin=-1, xmax=4, ymin=-0.05, ymax=1.15,
  xlabel={$x$}, ylabel={$F(x)$},
  width=9cm, height=5.5cm,
  xtick={0,1,2,3}, ytick={0,0.125,0.5,0.875,1},
  yticklabels={0,1/8,4/8,7/8,1},
  axis lines=middle,
]
\addplot[thick, blue] coordinates {(-1,0) (0,0)};
\addplot[thick, blue] coordinates {(0,0.125) (1,0.125)};
\addplot[thick, blue] coordinates {(1,0.5) (2,0.5)};
\addplot[thick, blue] coordinates {(2,0.875) (3,0.875)};
\addplot[thick, blue] coordinates {(3,1) (4,1)};
\addplot[only marks, mark=*, blue] coordinates {(0,0.125)(1,0.5)(2,0.875)(3,1)};
\addplot[only marks, mark=o, blue] coordinates {(0,0)(1,0.125)(2,0.5)(3,0.875)};
\end{axis}
\end{tikzpicture}
\end{center}
\end{frame}

\begin{frame}{Gráfico de $F(x)$ para una v.a. discreta (cont.)}
\begin{alertblock}{Nota}
El punto lleno indica que $F$ incluye el valor (continuidad a derecha); el punto vacío indica el límite por izquierda.
\end{alertblock}
\end{frame}

\begin{frame}{Gráfico de $F(x)$ para una v.a. continua}
\begin{example}
Si $X$ tiene f.d.a. $F(x) = 1-e^{-x}$ para $x\ge 0$ y $F(x)=0$ para $x<0$ (veremos que es la distribución \emph{exponencial}):
\end{example}
\begin{center}
\begin{tikzpicture}
\begin{axis}[
  xmin=-1, xmax=5, ymin=-0.05, ymax=1.15,
  xlabel={$x$}, ylabel={$F(x)$},
  width=9cm, height=5.5cm,
  axis lines=middle,
  domain=-1:5, samples=100,
]
\addplot[thick, red] coordinates {(-1,0)(0,0)};
\addplot[thick, red, domain=0:5] {1-exp(-x)};
\end{axis}
\end{tikzpicture}
\end{center}
Es una función continua y creciente, sin saltos: no hay valores puntuales con probabilidad positiva.
\end{frame}

\section{Ejemplo integrador}

\begin{frame}{Ejemplo resuelto}
\begin{example}
Sea $X$ una v.a. con f.d.a.
$$F(x) = \begin{cases} 0 & x<0 \\ x^2/4 & 0\le x <2 \\ 1 & x\ge 2\end{cases}$$
Calcular: (a) $P(X\le 1)$, (b) $P(0.5 < X \le 1.5)$, (c) $P(X>1)$.
\end{example}

\begin{block}{Resolución}
\begin{itemize}
\item[(a)] $P(X\le 1) = F(1) = 1^2/4 = 0.25$.
\item[(b)] $P(0.5<X\le 1.5) = F(1.5)-F(0.5) = \dfrac{1.5^2}{4} - \dfrac{0.5^2}{4} = \dfrac{2.25-0.25}{4} = \dfrac{2}{4}=0.5$.
\item[(c)] $P(X>1) = 1-F(1) = 1-0.25 = 0.75$.
\end{itemize}
Como $F$ es continua (no tiene saltos), $X$ es una variable aleatoria \textbf{continua}.
\end{block}
\end{frame}

\section{Resumen}

\begin{frame}{Resumen de la clase}
\begin{block}{Ideas clave}
\begin{itemize}
\item Una variable aleatoria $X:\Omega\to\mathbb{R}$ traduce resultados de un experimento en números.
\item La función de distribución $F(x)=P(X\le x)$ concentra toda la información probabilística de $X$.
\item Propiedades de $F$: no decreciente, continua a derecha, $F(-\infty)=0$, $F(+\infty)=1$.
\item $P(a<X\le b)=F(b)-F(a)$; $P(X>x)=1-F(x)$; $P(X=x)=F(x)-F(x^-)$.
\item $X$ discreta $\Leftrightarrow$ $F$ escalonada; $X$ continua $\Leftrightarrow$ $F$ continua (existe densidad $f$).
\end{itemize}
\end{block}

\begin{alertblock}{Próxima clase}
Estudiaremos en detalle las distribuciones discretas más importantes: Bernoulli, Binomial, Pascal e Hipergeométrica.
\end{alertblock}
\end{frame}

\end{document}
